% TRANSLATION UNIT: translate this file only.
% ===========================================================================
%  04 — Marco referencial
%  The \section title is the canonical (Spanish) chapter name; the body prose
%  is written in ENGLISH and translated to Spanish later.
% ===========================================================================

\section{Marco referencial}

% DONE: Theoretical frame written in English. Sources read and cited: Lusardi
% (2019), Amagir et al. (2017), FINRA Foundation (2022), Dare et al. (2020),
% Cannistrà et al. (2024); Mancone et al. (2024), Mishra et al. (2024),
% Al-Afeef & Alsmadi (2025), Kovács & Terták (2024), Menberu (2024),
% Koskelainen et al. (2023), Gumilar et al. (2024), Lotter & Okoro (2024);
% Ren et al. (2023), Miao et al. (2020), Dhandapani (2025), Lange et al. (2019),
% Richardson (n.d.). No source was cited whose body was not read.

% ---------------------------------------------------------------------------
% CLAIMS INVENTORY — every prose claim is listed here with its evidence.
% The prose below makes only these claims.
% ---------------------------------------------------------------------------
% EVIDENCE: Big Three = interest compounding, inflation, risk diversification →
%   /tmp/fuentes/lusardi2019financial.md §2.1.
% EVIDENCE: ~30 % answer all Big Three; ~50 % answer interest+inflation →
%   /tmp/fuentes/lusardi2019financial.md §2.2.
% EVIDENCE: one extra correct answer → 3–4 pp higher planning probability
%   (10 pp in the Netherlands) → /tmp/fuentes/lusardi2019financial.md §3.
% EVIDENCE: "one size does not fit all"; tailored, targeted, evaluated programs →
%   /tmp/fuentes/lusardi2019financial.md §4.
% EVIDENCE: PISA 10 % max proficiency; ~1 in 5 lack basic skills →
%   /tmp/fuentes/lusardi2019financial.md §4.
% EVIDENCE: school programs improve knowledge/attitudes; actual behavior
%   scarcely known; short-term; longitudinal needed → /tmp/fuentes/amagir2017review.pdf p.75.
% EVIDENCE: experiential learning / learning by doing is the promising pedagogy →
%   /tmp/fuentes/amagir2017review.pdf p.76.
% EVIDENCE: 76 RCTs, 33 countries, 160,000+ individuals; knowledge "large",
%   behavior "medium", low cost → /tmp/fuentes/finrafoundation2022financial.pdf p.1.
% EVIDENCE: not all programs effective; no evidence on decay of effects →
%   /tmp/fuentes/finrafoundation2022financial.pdf p.5–6.
% EVIDENCE: Solomon four-group, N=2,650 Dutch fifth grade; improved
%   "performing transactions" but not "responsible spending" →
%   /tmp/fuentes/dare2020effect.pdf p.93, p.100.
% EVIDENCE: game-based escape room RCT, 2,220 students, 4 countries,
%   +0.313 SD (0.182–0.618 by country) → /tmp/fuentes/cannistra2024impact.pdf p.825.
% EVIDENCE: only two studies compare game vs. traditional teaching; no
%   significant differences → /tmp/fuentes/cannistra2024impact.pdf p.827.
% EVIDENCE: effective programs integrate experiential learning, digital tools,
%   parental involvement; longitudinal studies needed →
%   /tmp/fuentes/mancone2024youth.pdf p.1, p.11.
% EVIDENCE: gamification mechanics motivate participation/retention; programs
%   without behavioral focus may not change habits → /tmp/fuentes/mancone2024youth.pdf p.9–11.
% EVIDENCE: five predictors explain ~71 % of variance in women's financial
%   decision-making (SEM, N=385) → /tmp/fuentes/mishra2024digital.pdf p.1.
% EVIDENCE: digital literacy → financial behavior/inclusion/well-being,
%   N=232, all four hypotheses supported → /tmp/fuentes/alafeef2025digital.md §Abstract.
% EVIDENCE: digital divide within generations; social media as channel;
%   education must teach information assessment → /tmp/fuentes/kovacs2024thematic.pdf p.491.
% EVIDENCE: technology-mediated education is scalable/cost-effective/reaches
%   underserved but faces infrastructure, access, language, exclusion barriers →
%   /tmp/fuentes/menberu2024technology.md §Abstract.
% EVIDENCE: three themes (fintech, digital financial behavior, behavioral
%   interventions); little known about digitalization's effect on behavior →
%   /tmp/fuentes/koskelainen2023financial.pdf p.507 (abstract).
% EVIDENCE: digital financial literacy is a precondition for digital financial
%   inclusion; coverage grows but usage barriers remain →
%   /tmp/fuentes/gumilar2024digital.pdf p.1563.
% EVIDENCE: higher-education initiatives vary in design and success; ongoing
%   evaluation needed → /tmp/fuentes/lotter2024review.pdf p.431.
% EVIDENCE: microservice architecture = independently deployable, loosely
%   coupled services around business capability; database per service; saga;
%   API gateway; trade-offs → /tmp/fuentes/richardsonndpattern.txt.
% EVIDENCE: educational big-data platform on microservices; Docker isolates
%   environments; components update independently → /tmp/fuentes/miao2020microservice.pdf p.1, p.11.
% EVIDENCE: microservices online learning platform; decentralization improves
%   stability/availability and lowers maintenance → /tmp/fuentes/ren2023design.pdf p.455–457.
% EVIDENCE: financial services monolith→microservices for real-time processing,
%   scalability, resilience; DDD for decomposition →
%   /tmp/fuentes/dhandapani2025microservices.pdf p.145.
% EVIDENCE: legacy insurance core monolith lacked flexibility/scalability/fault
%   tolerance; microservices with database-per-service + API gateway + Docker
%   guarantee them → /tmp/fuentes/lange2019microservices.pdf §5–6.

Financial literacy is low everywhere it has been measured, and the deficit is
not a curiosity of one country but a stable, cross-national finding. This
chapter assembles the evidence that justifies the two decisions that define
this project: to teach financial concepts through a digital platform, and to
build that platform as a set of microservices. The argument proceeds in three
strands---the evidence on financial literacy and what education can claim, the
evidence on digital and gamified delivery, and the architectural
literature---and closes at the point where the three converge on the platform
described in the remaining chapters.

\subsection{Financial literacy and what education can claim}

The most widely replicated measure of financial literacy is the set of three
questions that \textcite{lusardi2019financial} calls the \enquote{Big Three},
which test the capacity to calculate compound interest, the understanding of
inflation, and the understanding of risk diversification. Across the fifteen
countries of the FLat World comparison, the share of respondents who answer
all three correctly is around 30 \%, and only around 50 \% answer the two
questions on interest rates and inflation correctly; knowledge of risk
diversification is consistently the weakest of the three. The deficit is also
concentrated where it hurts most: fewer than 30 \% of Americans answer all
three questions correctly before age 40, and across countries women are less
likely than men to answer correctly and more likely to answer \enquote{do not
know}, a pattern Lusardi reads as awareness of a knowledge gap rather than
mere modesty. The deficit is not merely descriptive. Answering one additional question correctly is associated
with a 3--4 percentage point greater probability of planning for retirement,
and 10 percentage points in the Netherlands, where the association is
strongest \parencite{lusardi2019financial}. The stakes are not small: in a
life-cycle model, differences in financial knowledge account for 30 to 40 \%
of retirement wealth inequality in the United States
\parencite{lusardi2019financial}. Low literacy also correlates with
more expensive borrowing, weaker debt management, and less planning, which is
why Lusardi treats the \enquote{one size does not fit all} conclusion as the
design lesson: programs should be tailored to an identified audience, carry
clear objectives, and be evaluated rigorously. School-based exposure matters
because it reaches people before consequential decisions; in the PISA
assessments of 2012 and 2015 only 10 \% of 15-year-olds reached maximum
proficiency, and about one in five lacked even basic skills
\parencite{lusardi2019financial}.

What can education actually claim? \textcite{amagir2017review} reviews
school-based programs and finds that they reliably improve financial
knowledge and attitudes, while the evidence on actual behavior is scarce:
the studies measure mainly short-term effects, retention is small, and little
is known about what children and adolescents do with the knowledge afterward,
so the authors call for longitudinal experimental research. The same review adds
two cautions that matter for design: school programs can narrow the gender
gap, and self-reported knowledge and behavior may measure confidence and
intention rather than knowledge and action, so assessments of actual behavior
remain the weak point of the field. It identifies experiential learning---\enquote{learning by doing}, simulations,
and tasks tied to students' own life events---as the pedagogy most strongly
associated with knowledge gains. At the aggregate level, a meta-analysis
published by the \textcite{finrafoundation2022financial} synthesizes 76
randomized experiments across 33 countries and more than 160,000 individuals.
It finds positive effects on both knowledge and behavior, but the two effects
are not equal: the knowledge effect is judged \enquote{large} and the
behavioral effect \enquote{medium} under education benchmarks, at low average
cost. The same brief cautions that the vast variability across program designs
means that not all programs work, and that there is no evidence either way on
whether effects decay over time.

Two single-program evaluations make the limits of the \enquote{financial
education works} claim concrete. \textcite{dare2020effect} used a
quasi-experimental Solomon four-group design---a design chosen precisely so
that gains from repeated testing can be separated from gains from the
intervention---on a representative sample of 2,650 Dutch fifth-grade pupils.
The program improved knowledge and skills in performing transactions
effectively, but did not improve them in responsible spending. A program, in
other words, moves specific competencies, not literacy in general.
\textcite{cannistra2024impact} ran a randomized controlled trial of an online
game-based \enquote{Financial Escape Room} course with 2,220 students in four
countries, and found an average improvement of 0.313 standard deviations---but
the effect ranged from 0.182 in Estonia to 0.618 in Belgium; within countries, boys, native
speakers, and students on academic tracks benefited most, which the authors
attribute to differences in baseline literacy. The game required no teacher
mediation, which the authors read as evidence that a brief, low-cost,
fully digital game can move financial literacy at scale. Together these
strands support a measured conclusion: financial education is real but
modest and heterogeneous, it improves knowledge more reliably than behavior,
and its effectiveness depends on design---applied, experiential, and specific.

\subsection{Digital tools and gamification in financial education}

Digitalization changes not only how education is delivered but what literacy
itself must mean. \textcite{koskelainen2023financial} map the intersection of
finance and digitalization into three themes---fintech, financial behavior in
digital environments, and behavioral interventions---and propose a framework
for digital financial literacy, while conceding that very little is known
about how digitalization actually influences financial behavior. The stakes are
concrete: the same review observes that digital payment and online shopping
lower the barriers to consumption and can increase spending, so digital
financial literacy is as much about self-protection as about convenience.
\textcite{kovacs2024thematic} add a practical warning: there is a digital
divide even within a single generation, social media has become a dominant
transmission channel, and financial education must now prepare people to
assess whether information received online is true or biased. Kovács and Terták also
note a generational consequence: as the digitalization of finance accelerates,
the transmission of financial practice from parents loses relevance and social
media and formal education take its place, which raises the bar for what a
curriculum must teach. \textcite{gumilar2024digital}, reviewing 2020--2024 studies, frame digital
financial literacy as a precondition for digital financial inclusion, and note
that expanding coverage does not remove barriers on the usage side. The
delivery evidence points the same way. \textcite{menberu2024technology}
reviews technology-mediated financial education in developing countries and
finds it scalable, cost-effective, and able to reach underserved populations,
but bounded by infrastructure, unequal access, language and cultural
barriers, and the exclusion of marginalized communities. Menberu closes with three gaps
the field has not yet filled: evaluation of long-term effects, comparative
analysis of delivery approaches, and inclusive methodologies for the
populations the digital divide leaves out. \textcite{lotter2024review} finds, for higher education, a diverse landscape
of initiatives with varying degrees of success, and urges ongoing evaluation
rather than a single verdict.

The survey evidence for digital literacy is correlational, not causal, and it
should be read as such. \textcite{mishra2024digital} model the financial
decision-making of 385 Indian women with structural equation modeling and
report that five predictors---including digital financial literacy and
financial accessibility---explain around 71 \% of the variance in the
outcome. The same model finds that financial resilience moderates how
strongly decision-making translates into investment intentions, a reminder
that skills operate inside a context of resources and shocks.
\textcite{alafeef2025digital} reach a similar conclusion from a
social-media sample of 232 respondents, with all four hypotheses linking
digital literacy to financial behavior, inclusion, and well-being supported.
Both designs establish association in self-reported samples; neither
establishes that teaching digital skills causes better financial outcomes.

Gamification, the most visible claim in this literature, is the least
supported by hard evidence. \textcite{mancone2024youth} list gamification
mechanics---leaderboards, badges, and virtual rewards---among the components
that motivate participation and retention, but their review is narrative, it
recommends longitudinal studies precisely because long-term effects are
unmeasured, and it warns that programs focused on knowledge dissemination
without behavioral change may not produce lasting habits. The stronger
caution comes from \textcite{cannistra2024impact}, who note that only two
studies have directly compared game-based with traditional instruction and
that both found no significant difference in learning. The honest reading is
therefore not that games teach more, but that they engage and scale at low
cost---an argument about reach and participation, not about superior
pedagogy. This literature is thin and partly contradictory, and the chapter
does not smooth that over: digital delivery has firmer support than
gamification itself.

\subsection{Microservices as an architectural response}

The architectural justification for this project rests on a definition and on
two bodies of applied precedent. \textcite{richardsonndpattern} defines the
microservice architecture as a set of independently deployable, loosely
coupled services organized around business capability, each owning its own
database, fronted by an API gateway, and coordinated across services through
patterns such as the saga; the same catalog is explicit about the trade-offs,
chiefly distributed complexity and eventual consistency. The pattern is
defined as a set of forces in tension: small, simple, independently deployable
services on one side, and the cost of distributed operations---complex
interactions, eventual consistency, and runtime coupling---on the other, which
is why the catalog pairs the architecture with its mitigations rather than
presenting it as free of cost. That definition is the anchor, because it names exactly the properties this project needs: the
content that teaches and the computation that calculates evolve at different
rates and under different load, and only independent deployment lets each
change without disturbing the other.

The educational precedent is real. \textcite{miao2020microservice} describe
Qunxian, a microservice-based big-data platform for online educational
applications, and justify the choice by the need to update components
independently as the underlying data-science stack changes, using Docker to
isolate incompatible runtime environments. The platform shares data,
computing power, and infrastructure across a multiuser JupyterHub
environment, which is precisely the kind of shared yet partitioned workload
that a single monolith would struggle to serve. \textcite{ren2023design} build an
online learning platform on microservices and report that the decentralized
structure improves stability and availability while lowering maintenance
cost. The financial precedent matters equally, because this platform handles
money. \textcite{dhandapani2025microservices}, an argumentative
industry-oriented paper, holds that moving banking and payment systems from
monoliths to microservices delivers real-time processing, scalability, and
resilience, with domain-driven design as the decomposition discipline, and reports that leading adopters update 8 to
10 times more frequently than organizations on monolithic architectures---the
concrete form of the independent-deployment claim.
\textcite{lange2019microservices} provide the concrete case: a legacy
insurance core whose monolith suffered poor flexibility, scalability, and
fault tolerance was migrated to four independent services, each with its own
database and a common API gateway, and the authors conclude that the
resulting architecture guarantees the scalability and fault tolerance the
monolith could not. None of this is a survey of the catalog; it is the
specific evidence that the pattern, applied to education and to finance,
answers the requirements this project sets for itself.

\subsection{Where the three strands converge}

The three strands point at the same artifact. Financial literacy is low and
its consequences are costly, so the platform must teach; education improves
knowledge reliably but behavior only modestly, so the platform must teach by
doing---calculators that compute and quizzes that exercise, rather than text
alone. Digitalization redefines what a student must know, so the platform must
train digital financial skills, not only analog ones. And the microservices
pattern is what lets the two halves of that mission---content and
computation---be built, deployed, and corrected independently. The project is
the intersection: a financial-education platform whose calculators are data,
whose quizzes are practice, and whose services are independently deployable.
Read as a whole, the framework commits the project to three positions rather
than three slogans: to teach by application because that is where the
behavioral evidence points, to treat digital skill as part of the subject
matter rather than only as its delivery channel, and to accept the costs of
the microservices pattern---distributed operations and eventual
consistency---because independent deployment is the property that lets
teaching and calculation evolve at different speeds. Those commitments are
the criteria against which the architecture and the interface described in
the remaining chapters can be judged.
